% Imporditud: tools/import_eksperimendid.py
% Allikas: Eesti füüsikaolümpiaadi eksperimendi ülesannete
%   kogu (Rael Kalda, 2023), ülesanne Ü39, lahendus L39.
\setAuthor{EFO žürii}
\setRound{piirkonnavoor}
\setYear{1999}
\setNumber{G E1}
\setDifficulty{2}
\setTopic{dünaamika}
\setVahendid{suurem kõvakaaneline raamat, silindriline keha, joonlaud}
\setPunktid{11}
\setAllikas{Eksperimendi kogu Ü39, lahendus lk 51}
\setLahendusseis{toores_ilma_skeemita}

\prob{Hõõrdetegurid}
Hinnake kui palju erinevad hõõrdetegurid libisemisel ja veeremisel.

\hint

\solu
Katse: Raamatukaanest teha kaldpind ja lasta sealt alla veereda ja libiseda katsekeha. Kaane serva kõrgust mõõdetakse selles asendis, kus keha hakkas liikuma. Mõõtes kaane pikkuse raamatu seljani, saab arvutada kaldenurga.

\textit{Žürii hindamisskeem on kogumikus lk 51; siia ei ole seda ümber kirjutatud.}
\probend
